\documentclass[11pt]{amsart}
\usepackage{amsmath,amssymb,amsthm,geometry,hyperref,booktabs}
\geometry{margin=1in}

\newtheorem{theorem}{Theorem}[section]
\newtheorem{proposition}[theorem]{Proposition}
\newtheorem{lemma}[theorem]{Lemma}
\newtheorem{corollary}[theorem]{Corollary}
\theoremstyle{definition}
\newtheorem{definition}[theorem]{Definition}
\theoremstyle{remark}
\newtheorem{remark}[theorem]{Remark}

\newcommand{\Q}{\mathbb{Q}}
\newcommand{\Z}{\mathbb{Z}}
\newcommand{\F}{\mathbb{F}}
\newcommand{\Zp}{\mathbb{Z}_p}
\DeclareMathOperator{\supp}{supp}
\DeclareMathOperator{\efd}{efd}
\newcommand{\leg}[2]{\left(\frac{#1}{#2}\right)}
\newcommand{\num}{\operatorname{num}}
\hypersetup{colorlinks=true, linkcolor=blue, citecolor=blue, urlcolor=blue}

\title[A conditional six-quantifier definition of $\Z$ in $\Q$]
      {A conditional six-quantifier definition of $\Z$ inside $\Q$:\\
       model-field assembly, its walls, and a machine-verified finite program}
\author{h10q project}
\date{Working draft, 2026-08-19 --- not submitted}

\begin{document}
\begin{abstract}
We describe a constructive route to a universal definition of $\Z$ inside
$\Q$ using six quantifiers, conditional on one explicit hypothesis $H$
about the solubility of a family of tied conics. Every finitary component
of the construction is machine-verified on a grid of $353$ cells
$(w,z)$ with $w$ an odd prime below $100$: all $353$ cells carry a
verified witness, and all $293$ aligned classes carry a certified
emergent-free member. Two structural theorems delimit the search space:
a generalized wall showing that every $b$ coprime to the cell data fails
on every branch --- valid for composite squarefree $A$ as well as prime
$A$ --- and a characterization theorem stating that an aligned class
exists for a cell exactly when the numerator of $z$ has a prime divisor
congruent to $1$ modulo $4$. The residual obstruction is analysed by an
exact $p$-adic sieve model whose matched predictions agree with the
observed obstruction at member level on an identical $8{,}760$-member
sample (marginal ratio $0.9993$; hit-conditioned union residual
$-2.9$; marginal-preserving permutation null $z = -0.09$). The analytic input is then identified exactly: we prove that
\emph{Schinzel's Hypothesis~H for the pair $\{q_1+Nt,\;G(t)\}$ implies the
per-class hypothesis}, indeed infinitely often, so the conditional
statement rests on a named classical conjecture rather than a bespoke
one; the Schinzel hypotheses themselves are verified exactly on all
$293$ canonical pairs. The verified computational horizon extends
off-grid to fifteen further cells, including three that are provably
obstructed for the canonical protocol --- the obstruction being
$w \equiv 11, 19 \pmod{20}$, a reciprocity phenomenon we prove and then
break with a non-canonical class family. We claim no unconditional
record: the published record remains the refereed ten-quantifier
definition, and our statement is conditional on Schinzel's Hypothesis~H
together with the finite class-existence side.
\end{abstract}

\maketitle
\tableofcontents

\section{Introduction}

Hilbert's Tenth Problem over $\Q$ asks for an algorithm which, given a finite
system of polynomials with rational coefficients, decides whether the system
has a common rational zero.  The problem remains open in both directions.
One standard route to undecidability would be an existential definition of
$\Z$ in $\Q$, since such a definition would transport the undecidability over
$\Z$ to rational points.  The converse is not known: failure of that route
would not imply decidability over $\Q$.
% src: README.md:3-9; THEOREMS.md:176-181

There is already a sharp first obstruction.  Cornelissen--Zahidi proved that
$\Z$ cannot be the projection of the rational points of a curve over $\Q$;
in particular, no polynomial with one existential witness defines $\Z$ in
$\Q$ \cite{cz}.  Thus two witnesses, equivalently a surface as witness
variety, form the first open existential case.  This is also exactly where
the available topological and thin-set lower-bound methods stop.
% src: THEOREMS.md:10-28,157-168
A different elliptic-curve route has its own cliff.  If only finitely many
primes are withheld from an $S$-integer ring, the indices $n$ for which
$x(nP)$ is integral are a finite union of congruence classes, relative to the
chosen integral Weierstrass model.  When nonempty this set has positive
density and is archimedeanly equidistributed on a real component, rather than
discrete.  Consequently the prime-index mechanism used in elliptic
Diophantine models degenerates at the cofinite-$S$ boundary; this statement
does not rule out unrelated predicates.
% src: THEOREMS.md:36-85,100-118

Koenigsmann changed the quantitative landscape by proving that $\Z$ is
purely universally definable in $\Q$ \cite{koenigsmann}.
% src: NOTES.md:46-58
Complementation turns this into the study of the existential rank
$\operatorname{rk}^{\exists}_{\Q}(\Q\setminus\Z)$.  Daans--Dittmann--Fehm
then supplied the geometric interpretation: for a non-quantifier-free
Diophantine set $D$ over $\Q$,
\[
  \operatorname{rk}^{\exists}_{\Q}(D)
  =\operatorname{efd}_{\Q}(D)+1,
\]
and $\operatorname{efd}_{\Q}(D)\le d$ precisely when $D$ is the image of a
morphism all of whose fibres have dimension at most $d$ \cite{ddf}.  The
quantifier race is therefore a fibre-dimension problem.  In particular, a
six-universal-quantifier definition would amount to the bound
$\operatorname{efd}_{\Q}(\Q\setminus\Z)\le5$, not merely to more efficient
logical notation.
% src: THEOREMS.md:388-409; THEOREMS.md:582-591

\subsection{The record chain and its audit}

The refereed benchmark is Daans's ten-quantifier universal definition
\cite{daans}.  Our source-level audit reconstructs its budget before applying
any fusion: two quaternion parameters, three witnesses for the auxiliary
$\Phi$ condition, one splitting witness, and two three-variable trace--norm
certificates give twelve variables.  Two uses of the intersection theorem
of Daans--Dittmann--Fehm save one coordinate each, giving ten; pairwise
reassociation cannot improve that proof shape.  The same audit isolates the
real places to seek an improvement: stronger simultaneous fusion, a cheaper
trace--norm primitive, or a joint treatment of the nested valuation-ring
conditions.
% src: THEOREMS.md:122-161

Sun subsequently announced an unrefereed seven-quantifier construction
\cite{sun}.  It retains Daans's bridge but replaces the six-witness
intersection block by a finite disjunction of three-witness quadrics, and
then uses one further fusion.  Our A2 audit checked the exact quaternion
identities, the block's soundness, the residue-field combinatorics, and
explicit bridge certificates.  It also traced the local--global completeness
argument through parameterized Hensel lifting, weak and norm approximation,
and Hasse--Minkowski.  No substantive missing implication was found, but the
posted preprint has stale cross-references at the decisive step, and its
local--global chain was hand-traced rather than machine-verified.  We
therefore call seven an \emph{unrefereed announcement}; ten remains the
refereed record.
% src: THEOREMS.md:185-265

\textbf{PROVED.}  The audit also corrects one ingredient of the unrefereed announcement
\cite{sun}.  Its finite-field selection lemma assumes a residue field larger
than $25$.  Exhaustion first showed that the small fields do not fail; more
conceptually, we proved the exact identity
\[
 N(\tau)=\frac{\#E_{\tau^{-2}}(\F_q)}{4},\qquad
 E_{\tau^{-2}}:y^2=x(x-1)(x-\tau^{-2}).
\]
Hasse's bound makes this positive over every odd finite field admitting an
admissible $\tau$.  Hence the threshold can be deleted: $\F_3$ is exceptional
only because it has no $\tau\notin\{0,\pm1\}$, while the preprint's other,
data-dependent exceptional places remain untouched.
% src: THEOREMS.md:278-320,361-377

\subsection{Contributions and scope}

The contributions of this paper are deliberately separated by status.

\begin{enumerate}
\item \textbf{CONDITIONAL --- a six-quantifier architecture.}
We tie a witness in Sun's quadratic block to Daans's congruence parameter by
$a=1+2s$.  Soundness, denominator clearing, the exact target-place selector,
and the global norm obstruction are proved; completeness is the explicitly
named hypothesis $H$.  If $H$ holds, then $\Q\setminus\Z$ has existential
rank at most six, $\Z$ has a six-quantifier universal definition, and the
essential fibre dimension is at most five.
% src: THEOREMS.md:413-475,477-565,582-591

\item \textbf{PROVED --- wall theorems.}
We prove the two-adic source--target alignment wall, show that branch choice
cannot remove the reciprocity obstruction, and prove the generalized wall:
every admissible $b$ whose odd support is coprime to the cell data fails.
The only available door is controlled non-coprimality.
% src: THEOREMS.md:721-797,1214-1239,1275-1332,1415-1431

\item \textbf{PROVED --- complete class characterization.}
For the prime-$b$ family, on every branch, an aligned class exists exactly
when the numerator of $z$ has a prime divisor congruent to $1$ modulo $4$.
The necessity holds for every admissible rational $a$, not only when
$1+4a^2$ is prime; the sufficiency is constructive.
% src: THEOREMS.md:1241-1258,1415-1431

\item \textbf{PROVED --- exhaustive finite verification.}
On the declared grid, every one of the $353$ cells has a replayable soluble
witness, and every one of the $293$ aligned classes has a certified
emergent-free member.  These are bounded, per-row theorems under the
proven-primality engine, not an extrapolation to all primes.
% src: THEOREMS.md:1363-1382,1481-1511; README.md:223-231

\item \textbf{EVIDENCE --- a mechanistic sieve with matched validation.}
Exact root lifting turns each recurring small-prime obstruction into an
odd-valuation $p$-adic mass.  On the identical $8{,}760$-member sample, the
matched marginal ratio is $0.9993$, the hit-conditioned union residual is
$-2.9$, and a marginal-preserving permutation gives $z=-0.09$.  These data
validate the small-layer mechanism in the measured window; they do not
supply a tail bound or a uniform density theorem.
% src: THEOREMS.md:1574-1604; data/l17_matched.jsonl:1-3

\item \textbf{PROVED --- an audited Schinzel reduction.}
For all $293$ canonical linear/octic pairs, exact certificates verify
irreducibility of the degree-$8$ factor, fixed-divisor and pair-product gcd
conditions, and every local condition.  The linear Dirichlet side is a
theorem.  The simultaneous prime-value assertion for the octic factor is
not proved; it is precisely the remaining Schinzel input.
% src: README.md:39-40,223-231; data/l17_schinzel_audit.jsonl:1

\item \textbf{PROVED --- off-grid horizon extension.}
The same certificate machinery closes and replays the four cells at $u=-1$,
with $w=101$, $w=103$, $w=107$, and $w=109$.  They lie beyond the verified
grid horizon $w\le97$.  This is an extension by four exact instances, not a
uniform all-$w$ theorem.
% src: THEOREMS.md:1626-1644; README.md:223-231
\end{enumerate}

\paragraph{What is not claimed.}
We claim \emph{no new unconditional quantifier record}: Daans's ten remains
the refereed record, and Sun's seven remains an unrefereed announcement.
The project hypothesis $H$ is open.  In particular, prime values of the
irreducible degree-$8$ factor in the prescribed progressions are open; the
closest unconditional square-class results have degree at most two, while
the degree-three result is conditional on the elliptic-curve Parity
Conjecture \cite{krumm}.  The finite grid, the matched sieve statistics, and
the clean local audit do not prove that missing prime-value statement.
% src: THEOREMS.md:259-265,1646-1654; data/litscout_h10q.md:6-13,24-28

\section{Model-field architecture: selection, witness tying, and the open norm problem}

\subsection{The finite-field selector is exact}

\begin{theorem}[PROVED: L4, the small-field exhaustion]
Let $q$ be an odd prime power with $q\leq 25$, let $\chi$ be the quadratic
character of $\F_q$ with $\chi(0)=0$, and let
$\tau\in\F_q\setminus\{0,\pm1\}$.  There is an $a\in\F_q$ such that
\[
 \chi(1+4a^2)=-1,
 \qquad
 \chi\bigl(-(1+4a^2)\,(1-\tau^2(1+4a^2))\bigr)=1.
\]
% src: THEOREMS.md:278-289
For $q=3$ the assertion is vacuous, because there is no admissible $\tau$.
For $q\in\{5,7,9,11,13\}$ the minimum number of choices is $2$, and for
$q\in\{17,19,23,25\}$ it is $4$.
% src: THEOREMS.md:287-294
\end{theorem}

Thus the constant $25$ in the unrefereed character-sum proof of
\cite{sun} is not a genuine small-field obstruction: among odd residue
fields, only the field of order $3$ remains, and there the issue is
admissibility rather than failure of the selector.
% src: THEOREMS.md:300-304

\begin{theorem}[PROVED: L5, the Legendre identity]
For every odd prime power $q$, every
$\tau\in\F_q\setminus\{0,\pm1\}$, and $\lambda=\tau^{-2}$, let
$E_\lambda:y^2=x(x-1)(x-\lambda)$.  Then
\[
N(\tau) := \#\bigl\{a \in \mathbb{F}_q : \chi(1+4a^2) = -1,\
\chi\bigl(-(1+4a^2)(1-\tau^2(1+4a^2))\bigr) = 1\bigr\} \;=\; \frac{\#E_\lambda(\mathbb{F}_q)}{4}.
\]
% src: THEOREMS.md:308-315
Consequently
\[
 \#E_\lambda(\F_q)\geq(\sqrt q-1)^2>0,
 \qquad N(\tau)\geq(\sqrt q-1)^2/4>0,
\]
so the hypothesis $|k|>25$ can be deleted from the lemma.
% src: THEOREMS.md:316-320
\end{theorem}

The proof is an exact point count, not an asymptotic substitution.  With
$u_a=1+4a^2$ and $w_a=-u_a(1-\tau^2u_a)$, expand the corrected indicator;
the hyperbola identity $\sum_b\chi(b^2+c)=-1$ for $c\neq0$ evaluates the
quadratic pieces.  Fibre-counting the map $a\mapsto u_a$ then gives
\[
4N=q+1+\sum_u\chi\bigl(u(u-1)(u-\lambda)\bigr)
   =\#E_\lambda(\F_q).
\]
% src: THEOREMS.md:322-359
This also explains divisibility by $4$: the Legendre curve has its full
rational $2$-torsion.
% src: THEOREMS.md:366-369

\subsection{The tied model and its six-variable bookkeeping}

Set $K=\Q$, $S=\{2\}$, $\pi=2$, and $u=1$.  Pull the bridge back along
$z\mapsto z^3$ and put
$A=1+4a^2$, $B=2b$, $\delta_\tau=1-A\tau^2$, and
$c=h(a,b,z^3)=a^2z^6g(a,b)/(1-z^3-a^2z^6)$.
% src: THEOREMS.md:413-420
The original block is
\[
\Psi_\tau(a,b,c):\ \exists y,r,s\ [\delta_\tau(c^2 - Ay^2) - 16B(r^2 - As^2) = 16].
\]
% src: THEOREMS.md:419-420
The architectural step is to identify its witness $s$ with the congruence
parameter by imposing $a=1+2s$.  We use the finite branch menu
\[
 \tau\in\left\{0,\frac{2a}{A},
 \tau^\dagger=\frac{1+2a^2}{A}\right\}.
\]
After clearing the nonzero denominators, the resulting two-witness block is
\[
\begin{aligned}
\Theta_*(a,b,c,s):\quad \exists y,r\quad&
\bigl[c^2-Ay^2-16Br^2=16-16ABs^2\bigr]\\
&{}\vee
\bigl[c^2-Ay^2-16ABr^2=16A-16A^2Bs^2\bigr]\\
&{}\vee
\bigl[a^4(c^2-Ay^2)+4ABr^2=4A^2Bs^2-4A\bigr].
\end{aligned}
\]
The three equations correspond respectively to
$\tau=0$, $\tau=2a/A$, and $\tau^\dagger$; on the last branch
$\delta_{\tau^\dagger}=-4a^4/A$ and
$\alpha=(2a^2)^2$ is a square.  A finite disjunction is represented by the
product of these polynomials in the same witnesses $(y,r)$, so adding the
square branch costs no variable.
% src: data/l19_classexist.jsonl (square-branch definition); THEOREMS.md:1100-1121
The candidate model is therefore
\[
F(z):=\exists b,s\ \bigl[(1+2s,b)\in\Phi_1^{\{2\}}\ \wedge\
\Theta_*(1+2s,b,h(1+2s,b,z^3),s)\bigr].
\]
% src: THEOREMS.md:433-439
Here $\Phi_1^{\{2\}}$ enforces $v_2(s)\geq0$ and
$b\in\mathcal O_2^\times$.  Membership in it has existential rank at most
$3$; the tied block has rank $2$ in $(y,r)$.
% src: THEOREMS.md:453-456,470-473
The intersection theorem of \cite{ddf} replaces their conjunction by
$3+2-1=4$ witnesses, and projection of the two base coordinates $(b,s)$
gives
\[
 2+(3+2-1)=6.
\]
% src: THEOREMS.md:453-465
If an input set has rank zero, the elementary conjunction bound is no
larger.  The two branches share $(y,r)$ and are combined by multiplying
their equations; denominator clearing introduces no inverse witness because
$A\equiv5\pmod8$ at $2$ and $1-z^3-a^2z^6$ cannot vanish over $\Q$.
% src: THEOREMS.md:441-463

For use by the later local model, write $Z=z^3$,
$D_z=1-z^3-a^2z^6$, and
$N_g(b)=16a^4b^2-A(b-1)^4$.  The denominator-cleared value polynomial is
\[
P(b)=16D_z^2A^2b^4-\delta_\tau a^4Z^4N_g(b)^2-32A^3s^2D_z^2b^5 .
\]
% src: THEOREMS.md:729-731,975-977
Finally define the tied target value
\[
M_\tau=16-\delta_\tau c^2-16ABs^2.
\]
% src: THEOREMS.md:527-535

\textbf{PROVED: soundness.}\quad\textbf{CONDITIONAL ON SCHINZEL H:
completeness.}
Soundness, $F(z)\Rightarrow z\in\bigcup_{w\notin S}\mathfrak m_w$, is
proved.  The converse is open unconditionally; Sections~\ref{sec:record}
--\ref{sec:irreducibility} prove it assuming classical Schinzel~H.
% src: THEOREMS.md:437-439,468-475
% src: data/l22_elimination.jsonl; data/l18_schinzel_implies_h.jsonl
If that lemma holds, the union outside $S$ is $\exists_6$; adjoining the
finitely many omitted maximal ideals preserves that count, so
$\Q\setminus\Z$ is $\exists_6$, $\Z$ is $\forall_6$, and
$\operatorname{efd}_\Q(\Q\setminus\Z)\leq5$.
% src: THEOREMS.md:582-591

\subsection{Local structure and the exact global obstruction}

\textbf{PROVED (W0).}

The first criterion, W0, is uniform even in the residue field $\F_3$.
For every finite field $k$ of odd order,
\[
\sum_{a\in k}\chi(1+4a^2)=-1.
\]
% src: THEOREMS.md:477-484
Hence some $a$ makes $A$ a nonsquare, and the canonical selector is
\[
\tau_*(a)=
\begin{cases}
0,&\chi(-1)=-1,\\[2pt]
2a/A,&\chi(-1)=1.
\end{cases}
\]
% src: THEOREMS.md:484-492
It satisfies $\delta_{\tau_*}\neq0$ and
$\chi(-\delta_{\tau_*}A)=1$: the two values of
$-\delta_{\tau_*}A$ are $-A$ and $-1$.  Since $s\mapsto1+2s$ is a
bijection, tying loses no residue class and removes L5's $\F_3$ vacancy.
% src: THEOREMS.md:493-501

\textbf{PROVED (W1).}

For W1, let $w$ be odd with
$v_w(A)=v_w(\delta_\tau)=0$, $v_w(b)=1$, $v_w(z)\geq1$, and
$v_w(s)\geq0$.  The cube pullback gives either $c=0$ or
\[
v_w(c)=2v_w(a)+6v_w(z)-2\geq6v_w(z)-2\geq4.
\]
% src: THEOREMS.md:504-508
Then the tied conic
\[
-\delta_\tau A\,y^2-16B\,r^2=16-\delta_\tau c^2-16ABs^2
\]
% src: THEOREMS.md:508
is soluble over $\Q_w$ if and only if
$\chi_w(-\bar\delta_\tau\bar A)=1$.  The right side is a unit congruent
to $16$; its two subtracted terms have valuations at least $8$ and $1$.
Thus the target-place test is independent of the tied value of $s$.
% src: THEOREMS.md:509-525

\textbf{PROVED (W2).}

W2 identifies the remaining obstruction.  Put
$E_\tau=\Q(\sqrt{-\delta_\tau AB})$.  The conic is equivalent to
\[
N_{E_\tau/\mathbb Q}
\bigl(\delta_\tau A\,y+4r\sqrt{-\delta_\tau AB}\bigr)
=-\delta_\tau A\,M_\tau,
\qquad M_\tau=16-\delta_\tau c^2-16ABs^2.
\]
% src: THEOREMS.md:527-538
Its global solubility is exactly
\[
M_\tau=0\quad\text{or}\quad
-\delta_\tau A M_\tau\in N_{E_\tau/\mathbb Q}(E_\tau^\times).
\]
% src: THEOREMS.md:539-545
The zero alternative is genuine.  In the two branches the data are
\[
\begin{array}{c|c|c}
\tau& E_\tau&M_\tau\\ \hline
0&\mathbb Q(\sqrt{-AB})&16-c^2-16ABs^2\\
2a/A&\mathbb Q(\sqrt{-B})&16-c^2/A-16ABs^2.
\end{array}
\]
% src: THEOREMS.md:546-552

\subsection{Structural corollaries and the remaining lemma}

\textbf{PROVED (L7a).}
At an odd place with $v(\delta_\tau A)=v(B)=0$ and $M_\tau\neq0$, the tied
conic is soluble when $v(M_\tau)$ is even; when that valuation is odd, it
is soluble exactly when $\chi_v(-\delta_\tau AB)=1$.  If also $v(A)=0$,
the corresponding untied ternary form is universal over $\Q_v$.
% src: THEOREMS.md:602-620
Accordingly the full candidate obstruction set is
\[
\{2,\infty\}\ \cup\ \{v:\ v(M_\tau)\ \text{odd}\}\ \cup\ \{v:\ v(A)\ne0
\ \text{or}\ v(B)\ne0\}.
\]
% src: THEOREMS.md:622-645
Coefficient-bad places are genuine and cannot be replaced merely by the
support of $AB$, because valuations can cancel in that product.
% src: THEOREMS.md:622-643

\textbf{CONDITIONAL (L7b, conditional on L6 soundness).}
No fixed admissible $(s,b)$ and branch admits rational functions
$Y(z),R(z)\in\Q(z)$ satisfying
\[
-\delta_\tau A\,Y(z)^2-16B\,R(z)^2=M_\tau(z)
\]
identically.  This rules out a uniform rational section, not pointwise
coverage of all relevant rational fibres by some fixed pair.
% src: THEOREMS.md:647-671

\textbf{PROVED (L7c).}
For $M_\tau\neq0$, local solubility is equivalent to
\[
t_v:=\bigl(-16\delta_\tau AB,\ -\delta_\tau A M_\tau\bigr)_v=1,
\qquad \prod_v t_v=1.
\]
% src: THEOREMS.md:673-682
Thus a failing tied conic is obstructed at an even number of places, at
least two, and repairs must change places in pairs.
% src: THEOREMS.md:676-684

\textbf{CONDITIONAL ON SCHINZEL H / OPEN UNCONDITIONALLY (assembly).}
For every odd target $w$ and every $z$ with $v_w(z)\geq1$, one must choose
$(s,b)$ with $(1+2s,b)\in\Phi_1^{\{2\}}$,
$s\equiv(\bar a_w-1)/2\pmod{\mathfrak m_w}$, $v_w(b)=1$, and a retained
branch such that
$M_\tau=0$ or
$-\delta_\tau A M_\tau\in N_{E_\tau/\Q}(E_\tau^\times)$.
% src: THEOREMS.md:557-565
In this local formulation the free $b$-direction must steer
$E_\tau=\Q(\sqrt{-\delta_\tau AB})$ against the moving odd-valuation and
coefficient-bad places.  Later sections implement the required choice on
the fixed canonical branch: L22 supplies irreducibility, L20 supplies the
admissible class, and classical Schinzel~H supplies the member.  Thus the
weak/norm-approximation problem remains open only unconditionally, not as
an extra premise of the conditional theorem.
% src: data/l22_elimination.jsonl; data/l20_admissible.jsonl
% src: data/l18_schinzel_implies_h.jsonl
% src: THEOREMS.md:701-714

\textbf{EVIDENCE.}
Bounded witness-switching searches support this steering mechanism, but
are not a proof of assembly and are not used to upgrade the conditional
six-quantifier conclusion.
% src: THEOREMS.md:686-700

\section{Obstruction walls and reciprocity steering}
\label{sec:walls}

\subsection{The alignment wall}

Put
\[
 a=1+2s,\qquad A=1+4a^2,\qquad B=2b,\qquad
 P=1-ABs^2.
\]
For the two original tied branches, the source square classes are
$D_0=P$ and $D_1=AP$, whereas the target norm fields are
$E_0=\Q(\sqrt{-2Ab})$ and $E_1=\Q(\sqrt{-2b})$.  The exact value
identities show that the branch value $u_\tau$ is a norm from
$\Q(\sqrt{D_\tau})$, while solubility asks that $Au_\tau$ be a norm
from $E_\tau$.  In both branches the mismatch is measured by
\[
 D_\tau\operatorname{disc}(E_\tau)
 \equiv -2AbP=:\kappa\pmod{\Q^{\times 2}}.
\]
% src: THEOREMS.md:728-744

\begin{proposition}[L8: absorption and the parity wall]
\textbf{PROVED.}  At an odd place $v$, assume the branch value
$M_\tau$ is nonzero and each of
$v(\delta_\tau A)$, $v(B)$, and $v(M_\tau)$ is even.  Then the tied
conic has a $\Q_v$-point.  Consequently an odd obstruction can occur
only where $A$, $B$, or $M_\tau$ has odd valuation.  On the other hand,
for every admissible pair, $v_2(P)=0$ and $v_2(\kappa)=1$.
% src: THEOREMS.md:746-761
\end{proposition}

\begin{proof}
After scaling the two conic variables and dividing by the even power
of $v$ contributed by $M_\tau$, one obtains a nondegenerate binary
quadratic form with unit coefficients representing a unit.  Over the
residue field it represents every nonzero class: in the anisotropic
case it is the norm form of the quadratic extension, whose norm is
surjective.  Hensel lifting gives the local point.  This is the
absorption mechanism.
% src: THEOREMS.md:746-757

Admissibility gives $v_2(ABs^2)\geq 1$, hence $P$ is a $2$-adic unit.
Thus $-2AbP$ has odd $2$-adic valuation.  The target discriminants
$-2b$ and $-2Ab$ likewise have valuation one, so the target is always
a field and can never equal the source quadratic field.
% src: THEOREMS.md:759-771
\end{proof}

The two exact-matching devices are therefore empty on the admissible
locus.  Indeed, matching source and target would require
$\operatorname{disc}(E_\tau)\equiv D_\tau$, equivalently
$\kappa\equiv1$ modulo squares; splitting the target would require
$\operatorname{disc}(E_\tau)$ itself to be a square.  The first
condition contradicts $v_2(\kappa)=1$, and the second contradicts the
odd valuation of the target discriminant.  This rules out precisely
these two identities, not every possible fixed square-class relation:
for a general class $j$, one would still have to prove that an
odd-multiplicity value prime inert in $\Q(\sqrt j)$ must occur, and no
such theorem is asserted here.
% src: THEOREMS.md:774-797

\textbf{EVIDENCE.}  The frozen rescues do not exhibit a further
support-controlled identity: all ten have odd-multiplicity wild primes,
and all measured symbols at those primes are $+1$.  In the recorded
session sweep, thirteen of 161 admissible pairs succeeded, while every
observed obstruction set had even cardinality.  These observations are
alignment statistics only; they do not prove a uniform mechanism or an
assembly theorem.
% src: THEOREMS.md:799-808

\subsection{Reciprocity steering}

\begin{lemma}[L9 steering lemma]
\textbf{PROVED.}  For $x,d\in\Q^\times$, let
\[
 T=\{2,\infty\}\cup
 \{q\text{ odd}:v_q(x)\neq0\text{ or }v_q(d)\neq0\}.
\]
Then $\prod_{v\in T}(x,d)_v=1$.  Hence, if all symbols in $T$ except
possibly one are $+1$, the remaining symbol is also $+1$.
% src: THEOREMS.md:826-834
\end{lemma}

\begin{proof}
Hilbert reciprocity gives the product over all places.  Outside $T$,
both entries are odd-adic units and their Hilbert symbol is $+1$, so
only the displayed finite set remains.
% src: THEOREMS.md:828-834
\end{proof}

For the tied conic take
$x=\alpha M_\tau$, $d=\alpha B$, and
$\alpha=-\delta_\tau A$.  Once every controlled symbol is $+1$, a
single wild odd-multiplicity prime $q_0$ with $v_{q_0}(d)=0$ is not an
obstruction: its symbol is forced by the lemma.  Thus the useful
arithmetic target is not complete smoothness, but a controlled part
together with one proven prime-power cofactor.
% src: THEOREMS.md:836-853

\begin{proposition}[L9a: the prime-$b$ symbol]
\textbf{PROVED.}  Fix an admissible $a$, a cell $(w,z)$, and a branch
$\tau$.  Let $b=\varepsilon q_1$, where $q_1$ is odd prime, and impose separately
\[
 v_{q_1}(\delta_\tau)=v_{q_1}(A)=v_{q_1}(a)
 =v_{q_1}(z)=v_{q_1}(D_z)=0.
\]
Then
\[
 (x,d)_{q_1}=\left(\frac{A}{q_1}\right).
\]
In particular, the moving place is aligned by choosing a residue class
with $\left(\frac{A}{q_1}\right)=+1$.
% src: THEOREMS.md:887-903
\end{proposition}

\begin{proof}
Clearing denominators and reducing at $b\equiv0\pmod{q_1}$ gives
$N_g(0)=-A$.  The hypotheses make this a unit, the denominator
contributes valuation four, and $v_{q_1}(x)=-4$.  Its unit part is
$A$ modulo squares, while $v_{q_1}(d)=1$; the Hilbert-symbol parity
formula gives the stated Legendre symbol.  Quadratic reciprocity makes
it constant on suitable congruence classes.
% src: THEOREMS.md:893-903
\end{proof}

The corrected progression freezes only genuine controlled places.
With $S_0=\{2\}\cup\supp(\alpha)$, the modulus contains the required
$2$-adic and $A$-character moduli and the Taylor exponents at places of
$S_0$; it does not freeze the moving prime.  For a coprime prime
$q_1$ in the selected class,
\[
 \supp(d)=S_0\cup\{q_1\},
 \qquad \left(\frac{A}{q_1}\right)=+1.
\]
Thus every place in the support of $d$ is controlled, and only the
odd-multiplicity primes of the remaining cofactor of $x$ are wild.
% src: THEOREMS.md:915-932
% src: THEOREMS.md:953-957

The need for the corrected modulus is witnessed by the recorded cell
$(w,u)=(73,5)$.  With $s=0$, $b_0=-29$, $\tau=2/5$, the superseded
recipe gives
$N=689120$.  The prime $q=244637629=29+355N$ has the same prescribed
class, sign, and value of $(A\mid q)$, but at $59$ the symbol changes:
$v_{59}(x(b_0))=0$ with symbol $+1$, whereas
$v_{59}(x(-q))=1$ with symbol $-1$.
% src: THEOREMS.md:942-951

\textbf{CONDITIONAL.}  Turning such a progression into an infinite
assembly family still requires the named Schinzel--H/Bunyakovsky input
that simultaneously makes $q_1$ prime and leaves a single prime
cofactor of $x$.  Reciprocity removes the last symbol once that input
is supplied; it does not prove the prime-value assertion.
% src: THEOREMS.md:958-970

\subsection{The generalized coprimality wall}

\begin{theorem}[L12a]
\textbf{PROVED.}  Let $a$ be admissible and suppose $A=1+4a^2$ is
prime.  Let $\tau$ be any branch, assume $v_A(z)\leq0$, and suppose every odd-valuation place of an admissible rational
$b$ is coprime to
$\{2,z,D_z,a,\delta_\tau,A\}$.  If $b_{\mathrm{sf}}$ denotes the
squarefree part of its odd part, then
\[
 \prod_{p\in\{A\}\cup\operatorname{oddsupp}(b)}(x,d)_p=-1.
\]
Thus at least one of these places obstructs every branch.
% src: THEOREMS.md:1275-1287
\end{theorem}

\begin{proof}[Proof sketch]
At each odd place $p$ of $b$, whether $p$ occurs in the numerator or
denominator, the dominant-term calculation makes $v_p(x)$ even and
its unit part
\[
 u_x\equiv A\pmod{\Q_p^{\times2}}.
\]
Therefore $(x,d)_p=(A\mid p)=(p\mid A)$, and multiplication over these
places gives $(b_{\mathrm{sf}}\mid A)$.  At $A$, the branch-uniform
parity calculation gives
$(x,d)_A=-(b_{\mathrm{sf}}\mid A)$: the minus sign is
$(2\mid A)=-1$, since $A\equiv5\pmod8$, while the branch-dependent unit
of $\alpha$ is a residue.  The two Jacobi factors cancel and leave
$-1$.
% src: THEOREMS.md:1288-1311
\end{proof}

The only door in this argument is non-coprimality.  If $b$ shares a
prime $f$ with the cell data, the reduction to the constant term at
$f$ and hence $u_x\equiv A$ can fail; the telescoping proof then no
longer applies.  This is a door, not by itself a solubility claim.
% src: THEOREMS.md:1325-1339

\subsection{Composite squarefree wall}

\begin{theorem}[L13e]
\textbf{PROVED.}  Let $a$ be admissible, put $A=1+4a^2$, and let $K$
be the squarefree kernel of its rational square class; this includes
composite, non-squarefree, and rational $A$.  Assume $v_p(z)\leq0$ for
every odd $p\mid K$ and that every odd-valuation place of an admissible
rational $b$ is coprime to $\{2,z,D_z,a,\delta_\tau,A\}$.  Then, for
every odd $p\mid K$ and every odd $\ell$ occurring to odd order in $b$,
\[
 (x,d)_p=\left(\frac{2b}{p}\right),\qquad
 (x,d)_\ell=\left(\frac{A}{\ell}\right),
\]
and
\[
 \prod_{p\mid K}(x,d)_p
 \prod_{\ell\in\operatorname{oddsupp}(b)}(x,d)_\ell
 =\left(\frac{2}{K}\right)^{1+v_2(b)}.
\]
% src: THEOREMS.md:1415-1425
\end{theorem}

\begin{proof}[Proof sketch]
At an odd $p\mid K$, the middle term
$-\delta_\tau a^4Z^4N_g^2$ strictly dominates the other terms of the
value polynomial, and $v_p(N_g)=0$.  Hence $v_p(x)$ is odd and its unit
class contains the complementary factor $K/p$.  Every prime divisor of
$K$ is $1\pmod4$, because it divides a value $n^2+4m^2$ with the data
coprime.  In the Hilbert-symbol formula, the complementary factors from
the two entries cancel.  This yields $(2b\mid p)$ independently of the
branch.
% src: THEOREMS.md:1415-1422

Multiplying over $p\mid K$ and over the moving primes of $b$ produces
the cross terms $(M\mid K)(K\mid M)$.  Quadratic reciprocity cancels
them factor by factor because every divisor of $K$ is $1\pmod4$.
Moreover $K\equiv A\equiv5\pmod8$, so $(2\mid K)=-1$.  Admissibility is
exactly $v_2(b)=0$, for which the product is $-1$; at the first
non-admissible boundary $v_2(b)=1$, the product is $+1$ and this wall
dissolves.  Even-exponent primes of $A$ do not enter $K$.
% src: THEOREMS.md:1419-1425
\end{proof}

It follows that the necessity direction of L11h is a theorem for every
admissible $a$: escaping the coprime wall requires a prime $p\mid K$
with $v_p(z)>0$, and every such $p$ is congruent to $1$ modulo $4$ and
divides the numerator of $z$.  Hence reachability implies that the
numerator of $z$ has a prime divisor congruent to $1$ modulo $4$.
This conclusion does not supply the separate prime-value input needed
for assembly.
% src: THEOREMS.md:1426-1431

\section{Class Dichotomy, Characterization, and Escapes}

Fix an admissible parameter $a$, put
\[
 A=1+4a^2,\qquad \delta_\tau=1-A\tau^2,
 \qquad \alpha=-\delta_\tau A,
\]
and first consider the prime-member family $b=\varepsilon q$, with
$\varepsilon\in\{\pm1\}$ and $q$ prime.  An \emph{aligned class} means a
nonempty prime progression in which the Hilbert symbols at every controlled
finite place, at the moving place $q$, and at the real place are all $+1$.
This is deliberately a frozen-side notion: odd-valuation primes emerging from
the value of $P(b)$ are a separate obligation.
% src: THEOREMS.md:973-991

\subsection{The L10 dichotomy}

\begin{proposition}[Canonical alignment and the opposite wall; \textnormal{PROVED in the stated scope}]
On the canonical branch $\tau=2a/A$, the prime-$b$ family is aligned with the
canonical square class $-1$ in the exact sense that
\[
 \delta_\tau=\frac1A,\qquad \alpha=-1.
\]
On the branch $\tau=0$, if $a$ is odd, $A$ is prime,
$A\nmid 2zD_z$, and the prime $q\ne A$ satisfies the prescribed unit
conditions, then no aligned class exists.
\end{proposition}
% src: THEOREMS.md:1003-1014
% src: THEOREMS.md:1016-1042

The first assertion is an identity, not permission to discard the places
above $A$: although $\alpha$ is an $A$-unit, $\delta_\tau=1/A$ forces those
places back into the controlled set.  This distinction is essential; the
obsolete grid count obtained by omitting them is not used here.  For the
second assertion, $\tau=0$ gives $\delta_\tau=1$ and $\alpha=-A$.  The local
calculation at $A$ yields
\[
 (x,d)_A\,\leg{A}{q}=\leg{-2\varepsilon}{A}=-1.
\]
Indeed admissibility gives $A\equiv5\pmod 8$, so
$\leg{-1}{A}=+1$ and $\leg{2}{A}=-1$, independently of the sign
$\varepsilon$.  Two required symbols are therefore permanently
anti-correlated.  This is a theorem for prime integral $A$ under the displayed
hypotheses; it is not a claim extracted from a finite failure search.
% src: THEOREMS.md:1016-1042

\subsection{Free branches, an immovable wall}

Branch completion itself costs nothing.  If $\tau=p/q$ is rational in the
existing parameters and its denominator is proved nonzero on the base
locus, adjoining that branch introduces no inverse witness.  Admissibility
makes $A$ a nonsquare in $\Q_2$, hence $\delta_\tau\ne0$; the norm identity
in the unrefereed argument of Sun~\cite{sun} is uniform in $\tau$; and a
finite disjunction of branch equations is represented by their product in
the same witnesses.  Explicit choices meet all four square classes of
$\langle-1,A\rangle$, including the square branch
\[
 \tau^\dagger=\frac{1+2a^2}{1+4a^2},\qquad
 \alpha=(2a^2)^2.
\]
% src: THEOREMS.md:1100-1146
% src: RESULTS.md:89-91

The corrected controlled set for a prime member is
\[
 S=\{2,3,5,7\}\cup\supp(\alpha)\cup\supp(\delta_\tau).
\]
Outside $S$, the fixed content of $P$ has even valuation; after that content
is removed, the reduction is a nonzero polynomial whose degree prevents a
permanent odd valuation at the remaining primes.  Thus no fixed bad place is
silently omitted.
% src: THEOREMS.md:1189-1209

More importantly, changing branches does not change the obstruction.  For
prime $A$, admissible $a$, and $v_A(z)\le0$, every branch in the relevant
square-class group satisfies
\[
 (x,d)_A(x,d)_q=-1.
\]
Hence at least one controlled symbol is negative.  This branch-independent
wall was checked exactly on $2{,}136$ extended instances across all four
explicit branches; the computation confirms the proved identity rather than
replacing it.
% src: THEOREMS.md:1214-1239
% src: RESULTS.md:95

\begin{theorem}[The L11h characterization; \textnormal{PROVED}]
For the prime-$b$ family, an aligned class exists at the cell $(w,u)$, with
$z=wu$, if and only if the numerator of $z$ has a prime divisor
$p\equiv1\pmod4$.  On the finite grid this selects $190$ of $353$ cells; the
other $163$ cells are unreachable by this family on every branch.
\end{theorem}
% src: THEOREMS.md:1241-1265
% src: RESULTS.md:96

\begin{proof}[Proof sketch]
For necessity, the generalized wall says that alignment is impossible unless
some odd prime $p\mid A$ also has $v_p(z)>0$.  Write $a=m/n$ in lowest terms.
Such a prime divides $n^2+4m^2$, and therefore
$(2m n^{-1})^2\equiv-1\pmod p$; consequently $p\equiv1\pmod4$.  The
factor-by-factor telescoping argument applies also when $A$ is composite or
rational, so this necessity is \textsc{proved} for every admissible $a$, not
merely supported by a search.
% src: THEOREMS.md:1415-1431

Conversely, let $p\equiv1\pmod4$ divide the numerator of $z$.  Choose odd
representatives $m,n$ satisfying
$n^2\equiv-4m^2\pmod p$ and set $a=m/n$.  Then $a$ is admissible and
$v_p(A)>0$, so the factor shared by $A$ and $z$ breaks the wall.  Freezing the
complete controlled set gives an explicit aligned progression.  The
construction, including its escape prime, chosen $a$, branch, and modulus,
is replayed on all $190$ predicted rows.  This constructive sufficiency and
the preceding necessity make the equivalence the structural centre of the
class analysis.
% src: THEOREMS.md:1253-1259
\end{proof}

\subsection{Non-coprime escapes and complete grid coverage}

The characterization concerns $b=\varepsilon q$.  The only effective door
through the generalized coprime wall is to share a cell prime deliberately:
\[
 b=\varepsilon f q,\qquad f\mid z.
\]
This produces aligned escape classes at $103$ of the $163$ prime-family-walled
cells.  Their frozen-side alignment is \textsc{proved} row by row; the
remaining $60$ walled cells with no such escape are closed instead by direct,
fully soluble witnesses.
% src: RESULTS.md:98-100

The Schinzel audit also keeps six legacy residual rational/composite-$b$
witnesses separate as pointwise certificates.  They are unconditional local
witnesses, but they are not promoted into prime progressions and give no
prime-value theorem.  The eventual last cell, $(89,-2)$, was closed by the
replayed square-branch witness
$(a,b,\tau)=(3,89/367,19/37)$, completing cell coverage.
% src: README.md:40
% src: THEOREMS.md:1374-1382

\begin{center}
\begin{tabular}{lcc}
\hline
verified object & coverage & status\\
\hline
L11 prime-$b$ aligned classes & $190/190$ & \textsc{proved per row}\\ % src: RESULTS.md:96
non-coprime escape classes & $103/103$ & \textsc{proved per row}\\ % src: RESULTS.md:98-99
all aligned classes, with a zero-bad member & $293/293$ & \textsc{proved per row}\\ % src: RESULTS.md:108
all grid cells, with a soluble witness & $353/353$ & \textsc{proved per row}\\ % src: RESULTS.md:100
\hline
\end{tabular}
\end{center}
The class total is exactly the disjoint sum of the L11 and escape families;
the cell total additionally uses direct pointwise witnesses.  Neither total
asserts a uniform theorem beyond the verified grid.
% src: THEOREMS.md:1481-1509

\subsection{Why no square-class shortcut remains}

\begin{proposition}[L13c; \textnormal{PROVED per row}]
For every one of the $706$ cell--branch rows, $P$ is irreducible of degree
$8$ over $\Q$.  Consequently no constant-square or weak square-factor form
collapses the emergent-place obligation on the grid.
\end{proposition}
% src: THEOREMS.md:1384-1393
% src: RESULTS.md:101

The boundary between theorem and experiment is sharp.  The irreducibility
certificates are \textsc{proved}.  By contrast, the $137.5$ million exact
square-class tests with no admissible hit are \textsc{evidence} against a
hidden special-value shortcut, not a universal nonexistence theorem.
Likewise the observed Galois group $D_8\wr C_2$ is \textsc{evidence}: its
cycle types matched all $28{,}240$ certified unramified Frobenius samples,
but no group-identification proof is claimed.
% src: THEOREMS.md:1393-1402
% src: RESULTS.md:101

Finally, the earlier zero-bad search already supplied $24$ certified members
across $15$ distinct cells; those are \textsc{proved instances}, while the
rough quarter-density inferred from the sample is only \textsc{evidence}.
The later per-class replay strengthens the finite conclusion to the
$293/293$ line in the table, without changing the open uniform problem.
% src: THEOREMS.md:1404-1413
% src: THEOREMS.md:1481-1510

\section{The verified quantitative layer}
\label{sec:sieve}

The class construction reduces the residual local problem to a moving
arithmetic progression.  For an aligned class write
\[
 Q(k)=q_1+kN,\qquad b(k)=\varepsilon fQ(k),
\]
% src: THEOREMS.md:1529-1537
with the controlled Hilbert symbols already fixed to be positive.  A
member is called \emph{emergent-free} when the remaining odd-valuation
places contribute no negative symbol.  The purpose of the quantitative
layer is twofold: to certify this condition on the finite grid and to
separate the exact local mechanism from the statistical model suggested
by the resulting population.

\subsection{L14: instance verification}

\begin{proposition}[Finite instance of hypothesis $H$]
\textbf{PROVED (per row).}  All $293/293$ aligned classes in the
finite grid have a certified emergent-free member; these comprise $190$
L11 classes and $103$ escape classes.  Every stored row was replayed
through the exact cofactor ladder and returned verdict \emph{zero}.
% src: THEOREMS.md:1481-1486
\end{proposition}

The proposition is deliberately finite.  Its acceptance condition is
an aligned class together with an exact zero verdict from the member
ladder.  A tied-conic or ramification computation may also be attempted
as an auxiliary cross-check, but a budget refusal there is accepted and
logged: it neither invalidates a ladder certificate nor supplies
positive evidence.  In particular, no refusal is promoted to a
primality or solubility assertion.  This distinction is essential
because the general, all-cell form of hypothesis $H$ remains open even
though the displayed grid instance is complete.
% src: THEOREMS.md:1485-1487
% src: THEOREMS.md:1503-1511

\subsection{L15: the exact remainder and counting laws}

For a certified member, let $R$ be the remainder of $P(b)$ after the
frozen support, the fixed shell, and the smooth-emergent primes have
been stripped.  An independent replay gives the following exhaustive
shape table.

\begin{table}[ht]
\centering
\begin{tabular}{@{}lcll@{}}
\toprule
Ladder shape & Rows & Exact form of $R$ & Size data \\
\midrule
proved-prime & $197$ & one prime, never $pC^2$ & median $48$ digits; maximum $119$ \\ % src: THEOREMS.md:1522-1523
factorint & $96$ & squarefree product, $e\in\{2,3,4\}$ & multiplicities $76/15/5$ \\ % src: THEOREMS.md:1523-1525
\bottomrule
\end{tabular}
\caption{\textbf{PROVED (finite audit).}  Shapes of the stripped
remainder on all $293$ closure rows.}
% src: THEOREMS.md:1513-1525
\end{table}

Thus $R$ is squarefree on every audited row and $R=1$ never occurs.
% src: THEOREMS.md:1518-1524
The prime rung is not a probable-prime shortcut: its singleton cofactor
is accepted only after a proof of primality.  On a factorint rung the
integer $e$ counts distinct emergent factors, and the exact
factorization allows every corresponding Hilbert symbol to be checked.
The two shapes therefore reflect two proof routes through the same
local criterion, not two definitions of closure.

There is also an exact finite counting statement.  The first certified
member occurs at $k=0$ in $207/293$ classes, or $70.65\%$.
% src: THEOREMS.md:1529-1532
This is \textbf{PROVED} as a statement about deterministic playback of
the stored progressions.  By contrast, interpreting the searches as
samples from a stable member process is \textbf{EVIDENCE}: the measured
per-prime-member emergent-free rate is approximately $0.040$, and the
observed waiting time is approximately $15.7$ times the pure
Bateman--Horn \cite{batemanhorn} first-prime prediction for the same
progressions.
% src: THEOREMS.md:1533-1537
The exact histogram may therefore be used without qualification inside
the finite audit, while the rate and the comparison are empirical laws,
not uniform density theorems.

\subsection{L16: why the small layer has sieve shape}

Fix an emergent odd prime $p$.  On an aligned class, $b(k)$ is affine in
$k$, and the sign at $p$ is
\[
 (x_0,d_0)_p=\leg{2\alpha b(k)}{p}.
\]
On every checked recurring-prime row, this sign is determined by
$k\bmod p$.  The obstruction at $p$ is consequently not the whole root
set of $P(b(k))$ modulo $p$: it is the union of those root classes having
the bad signature, further restricted to odd valuation,
\[
 \mathcal B_p=\{k:v_p(P(b(k)))\ \text{is odd}\}
 \cap\{k\bmod p:\leg{2\alpha b(k)}{p}=-1\}.
\]
% src: THEOREMS.md:1546-1557

For a simple bad root, Hensel lifting gives an exact Haar mass
\[
 \sum_{j\geq0}\frac{p-1}{p^{2j+2}}=\frac{1}{p+1}.
\]
% src: THEOREMS.md:1555-1559
This is the valuation-parity correction to a naive exclusion sieve:
one must retain the odd levels of the Hensel tower rather than merely
delete a residue class modulo $p$.  A singular root has no such
one-line factor.  Its descendants are lifted recursively until the odd
and even masses are certified with a residual bound; if that cannot be
completed, the corresponding local product is explicitly labeled as
an exclusion or partial result.
% src: THEOREMS.md:1558-1562
Accordingly, the residue and valuation checks are \textbf{PROVED per
row}; the claim that this small layer supplies the useful global
layering model is \textbf{EVIDENCE}.  No independence theorem is hidden
in the product notation.
% src: THEOREMS.md:1541-1544

\subsection{L17: matched validation}

The full-grid residue test strengthens the preliminary structural
probe.  Across all $293$ classes and the window $k\in[0,119]$, it found
zero residue-determination counterexamples on the identical
$8{,}760$-member sample used for evaluation.
% src: THEOREMS.md:1577-1580
% src: data/l17_sieve.jsonl:1; data/l17_matched.jsonl:3
This is \textbf{PROVED} as a finite statement.  The predictive use of
the local masses is \textbf{EVIDENCE}; it was assessed only after the
predictions and observations had been put on exactly the same member
support.

\begin{table}[ht]
\centering
\begin{tabular}{@{}lrrl@{}}
\toprule
Matched statistic & Model & Observed & Diagnostic \\
\midrule
marginal odd hits & $4481.3$ & $4478$ & ratio $0.9993$ \\ % src: data/l17_matched.jsonl:3
hit-conditioned union & $3477.9$ & $3475$ & residual $-2.9$ \\ % src: data/l17_matched.jsonl:3
permutation null & --- & $3475$ & $z=-0.09$; tails $0.61/0.53$ \\ % src: data/l17_matched.jsonl:3
\bottomrule
\end{tabular}
\caption{\textbf{EVIDENCE.}  Matched small-layer evaluation on the same
$8{,}760$ members used on both sides.}
% src: data/l17_matched.jsonl:3
\end{table}

The marginal statistic counts odd bad-root hits, so a member can
contribute more than once; the union statistic asks whether a member is
hit at least once.  Their simultaneous agreement, together with the
marginal-preserving permutation null, shows no detectable overlap or
dependence discrepancy on this window.  It does not prove independence
outside the evaluated support.

A superficially different class-weighted comparison has estimand about
$+0.0195$, at $3.63$ standard deviations.
% src: data/l17_matched.jsonl:3
That quantity compares Haar-uniform local masses in $\Z_p$ with the
finite $k\in[0,119]$ window, rather than comparing predictions and
observations member by member on matched support.
% src: data/l17_matched.jsonl:3
It is therefore identified as the Haar--window gap, not as a failure of
the hit-conditioned model.

\subsection{L17 cofactor layer: reciprocity and its limit}

The large-cofactor layer has a narrower theorem than the small layer.
\textbf{PROVED:} on the prime rungs, reciprocity forces the singleton
cofactor sign to be $+1$ in all $197/197$ cases; on the complete closure
set, the cofactor-sign product is $+1$ on $293/293$ rows.
% src: THEOREMS.md:1605-1611
These are exactly the singleton and parity-product statements.  The
individual positive signs in factorint closure rows are selected by the
closure condition, so their residue tables are descriptive data and
cannot support an unconditional sign-distribution claim.
% src: THEOREMS.md:1611-1615

The attempted nonclosure measurement remains \textbf{INCONCLUSIVE}:
only $11/400$ sampled rows were resolved, and the predeclared power gate
failed.
% src: NOTES.md:835-845
The unresolved rows are refusals, not observations of either sign.
Consequently the cofactor layer is proved at reciprocity level and no
further; its nonclosure sign distribution remains measurement-limited.
This is also why the matched small-layer success must not be silently
upgraded into a complete rate theorem for emergent-free members.

\section{The conditional record and the Schinzel audit}
\label{sec:record}

\subsection{The statement, with its hypothesis exposed}

The member property below records the output required by L6; it is no
longer an additional conjectural premise.

\begin{definition}[Member property $\mathcal M$]
For every odd prime $w$ and every rational $z$ with $v_w(z)\geq1$, there
is a verified aligned class for $(w,z)$ and a prime member
$Q\equiv q\pmod N$ outside its finite excluded set whose bad emergent set
is empty.  Explicitly, for every odd prime
$\ell\notin S\cup\{Q\}$ with
$v_\ell(P(\varepsilon fQ))$ odd, the local Hilbert symbol
$(x_0,d_0)_\ell$ is $+1$, where $S$ is the frozen set of the class.
\end{definition}
% src: data/l19_classexist.jsonl (uniform class component)
% src: data/l18_schinzel_implies_h.jsonl (member contract)

\begin{theorem}[CONDITIONAL ON CLASSICAL SCHINZEL H: Theorem C]
Assume classical Schinzel's Hypothesis~H.  For every odd cell, use the
fixed canonical branch
\[
 \tau^\dagger=\frac{1+2a^2}{1+4a^2},
 \qquad \delta_{\tau^\dagger}=-\frac{4a^4}{1+4a^2}.
\]
Then the member property $\mathcal M$ holds and
$\Q\setminus\Z$ is Diophantine in $\Q$ with six unknowns,
realized by the L6 formula
\[
 F(z):=\exists b,s\;\bigl[(1+2s,b)\in\Phi_1^{\{2\}}\ \wedge\
 \Theta_*\bigl(1+2s,b,h(1+2s,b,z^3),s\bigr)\bigr],
\]
where, with $A=1+4a^2$, $B=2b$, and
$\tau^\dagger=(1+2a^2)/A$,
\[
\begin{aligned}
 \Theta_*(a,b,c,s):=\exists y,r\quad&
 [c^2-Ay^2-16Br^2=16-16ABs^2]\\
 &{}\vee[c^2-Ay^2-16ABr^2=16A-16A^2Bs^2]\\
 &{}\vee[a^4(c^2-Ay^2)+4ABr^2=4A^2Bs^2-4A].
\end{aligned}
\]
The three disjuncts remain the cleared equations for
$\tau=0$, $\tau=2a/A$, and $\tau^\dagger$, respectively, in the same
witnesses $(y,r)$.  They preserve the finite soundness menu and the
six-count; completeness in this theorem uses only the third, fixed
canonical disjunct.
% src: data/l19_classexist.jsonl (square-branch definition)
% src: data/l22_elimination.jsonl (fixed-Z-theorem row)
% src: CONDITIONAL.md:18-31
Consequently $\Z$ is definable in $\Q$ by a formula with six universal
quantifiers,
$\Z=\Q\setminus(F\vee\mathfrak m_2)$, and
$\operatorname{efd}_{\Q}(\Q\setminus\Z)\leq 5$.
% src: CONDITIONAL.md:14-25
\end{theorem}

This is a conditional definition, not an unconditional solution of
Hilbert's Tenth Problem over $\Q$.  Soundness of the finite branch menu
consumes no hypothesis.  For completeness, fixed-$Z$ vertical
irreducibility and Hilbert selection give a good canonical-branch $a$ in
the existing admissible progression; L20 supplies the class and all
Schinzel-pair conditions; Theorem~\ref{thm:schinzel} turns classical
Schinzel~H into an emergent-free member; and L6 gives the displayed
definition.  No per-cell irreducibility certificate is assumed.
% src: data/l22_elimination.jsonl (fixed-Z-theorem row)
% src: data/l20_admissible.jsonl; data/l18_schinzel_implies_h.jsonl

The quantifier saving is attributed precisely to Nicolas Daans, Philip
Dittmann, and Arno Fehm, \emph{Existential rank and essential dimension of
diophantine sets}, arXiv:2102.06941v5, Theorem~1.4 (equivalently
Corollary~5.11)~\cite{ddf}.  Their intersection theorem turns the ranks
three and two on the common base into $3+2-1=4$ witnesses; the two base
coordinates then give six unknowns.  The recorded rank--essential-dimension
identity gives the bound five.  These counts remain \textbf{CONDITIONAL}
because the set equality to $\Q\setminus\Z$ uses classical Schinzel~H.
% src: CONDITIONAL.md:31-48
% src: data/litscout_h10q.md:30-34

\subsection{What is proved, and what Schinzel must supply}

\textbf{PROVED (uniform class side and linear side).}
Theorem~\ref{thm:classexist} constructs an aligned class for every cell and
freezes all controlled Hilbert symbols while restricting the moving prime to
one reduced progression
\[
 q_1(t)=q_{1,0}+Nt,\qquad \gcd(q_{1,0},N)=1.
\]
The residue construction gives infinitely many proven-prime choices of the
base $q_{1,0}$, and Dirichlet's theorem supplies infinitely many prime values
of the progression.  This is the linear, unconditional half of the
member-existence problem; no probable-prime test or refusal is used as
evidence.
% src: data/l19_classexist.jsonl (meta, construction rows, and summary)

\textbf{CONDITIONAL (the sole conjectural cofactor side).}  Put
\[
 F_{\mathrm{cell}}(t)=P\bigl(\varepsilon f q_1(t)\bigr),
\]
using the primitive integer representative after removal of the fixed
rational scale recorded in the audit.  The L22 theorem and quantitative
Hilbert irreducibility make $F_{\mathrm{cell}}$ irreducible for a selected
admissible $a$.  Requiring it to be prime simultaneously with
$q_1(t)$ is the two-polynomial instance of classical Schinzel's
Hypothesis~H \cite{schinzel}.  The polynomial has degree eight.  A prime
cofactor gives an emergent-free member because the class has frozen the
controlled places and reciprocity forces the remaining single symbol.
After using the progression prime as variable, the pair is
$\{q_1+Nt,F_{\mathrm{cell}}(t)\}$.
% src: data/l22_elimination.jsonl (fixed-Z-theorem row)
% src: data/l20_admissible.jsonl; data/l18_schinzel_implies_h.jsonl
% src: THEOREMS.md:953-971
% src: CONDITIONAL.md:175-185
% src: data/l17_schinzel_audit.jsonl:1-294

\subsection{PROVED finite audit of the Schinzel pair}

The exact audit contains $293$ canonical records, split into $190$ L11
classes and $103$ escape classes.  The fields
\texttt{F\_irreducible}, \texttt{F\_value\_gcd}, and
\texttt{pair\_product\_gcd} respectively certify irreducibility on
$293/293$ rows, value gcd one on $293/293$ rows, and pair-product gcd one
on $293/293$ rows.  Every irreducibility certificate is a degree-eight
Frobenius certificate modulo a certified prime; both gcds were computed
exactly on $0\leq t\leq 2000$.  The pair-product computation is already a
certificate that no prime divides the product for every integer $t$: any
such fixed divisor would divide the computed gcd.  Thus all canonical
pairs pass the fixed-divisor and irreducibility conditions.  This is
\textbf{PROVED} for the stored pairs, not a prime-value theorem.

The finite audit does not by itself extrapolate to arbitrary rational
cells.  The L20 audit proves normalization, coprimality, absence of a fixed
divisor, frozen-unit and symbol constancy, and positive leading coefficient
uniformly.  Theorem~\ref{thm:fiber-irred} now supplies the missing uniform
input directly: for every fixed nonzero rational $Z$, $P(a,Z)$ is
irreducible in $\Q(a)[b]$.  Corollary~\ref{cor:admissible-irred} then
selects a good $a$ inside the character progression.  Thus the
$353$ recorded target certificates remain exact finite corroboration, but
they are no longer a premise for a new cell.
% src: data/l20_admissible.jsonl; data/l21_irred_generic.jsonl
% src: data/l22_elimination.jsonl (fixed-Z-theorem row)
% src: data/l17_schinzel_audit.jsonl:1-294

For a prime $p$ dividing $N$, the corrected residue count is
\[
 \nu_p=\#\{t\in\F_p:q_1(t)\not\equiv0\pmod p,
                  \ F_{\mathrm{cell}}(t)\not\equiv0\pmod p\}.
\]
The first condition is essential.  Since $p\mid N$ and
$\gcd(q_{1,0},N)=1$, the linear polynomial is the nonzero constant
$q_{1,0}$ modulo $p$; it must not be replaced by the parameter $t$.
Consequently the count is $p$, rather than $p-1$, whenever the cofactor has
no root modulo $p$.  The audit enumerated every $p\leq1000$ dividing the
row's modulus and obtained exactly this full count in every occurrence.
% src: data/l17_schinzel_audit.jsonl:1-294

\begin{table}[ht]
\centering
\small
\begin{tabular}{c|c|c|c}
primes $p$ & records containing $p\mid N$ & range of $\nu_p$ & mean density $\nu_p/p$\\ \hline
$2,3,5,7$ & $293$ each & $p$ & $1$\\
$11,37,89,97$ & $13$ each & $p$ & $1$\\
$13,53$ & $26$ each & $p$ & $1$\\
$17$ & $39$ & $p$ & $1$\\
$23$ & $38$ & $p$ & $1$\\
$29,41,61,73$ & $28$ each & $p$ & $1$\\
$43,47,67,83$ & $10$ each & $p$ & $1$
\end{tabular}
\caption{Corrected admissible-residue counts for primes dividing the progression modulus.}
% src: data/l17_schinzel_audit.jsonl:1
\end{table}

These local counts neither assume independence nor estimate a density of
simultaneous prime values.  Their role is narrower and exact: together with
the pair-product gcd, they rule out the elementary local obstruction for
each stored Schinzel pair.
% src: data/l17_schinzel_audit.jsonl:1

\subsection{Residual certificates and the literature boundary}

The six residual rows are deliberately segregated.  They are rational or
composite-$b$ witnesses verified by exact branch reconstruction and tied
status; they are not coerced into the prime family and are not evidence for
Schinzel's hypothesis.

\begin{table}[ht]
\centering
\small
\begin{tabular}{c|c|c|c}
cell $(w,z)$ & $(a,b)$ & $\tau$ & status\\ \hline
$(67,2)$ & $(1,4757)$ & $3/5$ & PROVED pointwise\\
% src: data/l17_schinzel_audit.jsonl:295
$(41,1/7)$ & $(25,4985)$ & $1251/2501$ & PROVED pointwise\\
% src: data/l17_schinzel_audit.jsonl:296
$(61,-2/7)$ & $(25,55571)$ & $1251/2501$ & PROVED pointwise\\
% src: data/l17_schinzel_audit.jsonl:297
$(61,2)$ & $(25,-71431)$ & $1251/2501$ & PROVED pointwise\\
% src: data/l17_schinzel_audit.jsonl:298
$(53,1/3)$ & $(-15,9)$ & $451/901$ & PROVED pointwise\\
% src: data/l17_schinzel_audit.jsonl:299
$(29,-2)$ & $(-15,-29/5)$ & $4055/2703$ & PROVED pointwise
% src: data/l17_schinzel_audit.jsonl:300
\end{tabular}
\caption{The six residual witnesses: unconditional certificates for their individual cells only.}
\end{table}

\textbf{OPEN (the unconditional published boundary).}  The exact
square-class reference is David Krumm, \emph{Squarefree parts of polynomial
values}, Journal de Th\'eorie des Nombres de Bordeaux \textbf{28} (2016),
no.~3, 699--724, arXiv:1407.4890~\cite{krumm}.  Its Theorem~1.3 and
Propositions~3.4--3.5 give the unconditional degree-one and degree-two
cases; Proposition~3.8 makes the degree-three case conditional on the
Parity Conjecture for elliptic curves over $\Q$.
% src: data/litscout_h10q.md:6-13

At prime arguments, the nearest unconditional results reach the cubic
boundary without prescribing a square class: Helfgott proves square-free
values of $f(p)$ for separable cubic $f$ \cite{helfgott}, while Reuss
proves $(d-1)$-free values rather than prime or prescribed-square-class
values \cite{reuss}.  Neither yields the required degree-eight
conclusion.
% src: data/litscout_h10q.md:15-28

The conic-bundle literature does not remove the gap.  The unconditional
Harpaz--Skorobogatov--Wittenberg theorem requires abelian constant fields
in the bad fibres \cite[Theorem~3.1 and Corollary~3.4]{hsw}.  The general
Schinzel reference in \emph{J. reine angew. Math.} \textbf{495} (1998),
1--28 is by Colliot-Th\'el\`ene, Skorobogatov and Swinnerton-Dyer
\cite{ctssd}, not by Colliot-Th\'el\`ene and Sansuc; the latter pair's
Schinzel paper is \cite{ctsansuc}.  No cited unconditional theorem supplies
simultaneous prime values for the present linear--octic pair.  Classical
Schinzel~H is therefore the sole conjectural input to the six-quantifier
theorem, which remains \textbf{CONDITIONAL}; no unconditional record or
solution of H10($\Q$) is claimed.
% src: data/litscout_h10q.md:36-39


\section{Schinzel's Hypothesis H: the fixed-class implication}
\label{sec:schinzel}

Theorem~\ref{thm:classexist} supplies the canonical aligned-class
construction for every cell.  Section~\ref{sec:irreducibility} proves, for
every fixed nonzero rational $Z$, that its octic is irreducible in
$\Q(a)[b]$, and quantitative Hilbert irreducibility selects an irreducible
specialization inside the existing character-admissible $a$ progression.
The L20 admissibility package is therefore unconditional.  For any such
fixed class, classical Schinzel's Hypothesis~H \cite{schinzel} is the sole
remaining input and supplies the member.
% src: data/l22_elimination.jsonl (fixed-Z-theorem row)
% src: data/l20_admissible.jsonl; data/l18_schinzel_implies_h.jsonl

\subsection{The operational contract}

A place is recorded as an obstruction only when its valuation is odd
\emph{and} its Hilbert symbol is $-1$: the kernel appends a place under
exactly that conjunction, and the ladder returns \emph{zero} when every
remaining place is proved with symbol $+1$.
% src: l13_filter.py smooth_emergent / cofactor_decide
Emergent-freeness therefore means \emph{no place of symbol $-1$}, not
\emph{no odd valuation outside the frozen set}. This distinction is
load-bearing: $197$ of the $293$ grid closures have exactly one
odd-valuation outside prime whose symbol is forced $+1$ by reciprocity.
% src: THEOREMS.md L14/L17 (197 prime-rung closures, 197/197 singleton +1)
The stronger reading would not follow from prime values, since a prime
value deliberately creates one odd place.

\subsection{The theorem}

Fix a verified aligned class with data $(a,\varepsilon,f,q_1,N)$, put
$Q(t) = q_1 + Nt$, $b(t) = \varepsilon f Q(t)$, and let
$S = \{2,3,5,7\} \cup \supp(\alpha) \cup \supp(\delta) \cup \supp(f)$.
Write the composed class polynomial as
\[
   F(t) \;=\; P\bigl(\varepsilon f (q_1 + Nt)\bigr) \;=\; c\,G(t),
\]
with $c > 0$ and $G \in \Z[t]$ primitive with positive leading
coefficient. The relevant hypothesis on $c$ is that its \emph{square
class} is supported on $S$. Literal $S$-support is false in general: for
$(w,z) = (3, 3/11)$, $a = 1$, $q_1 = 19$ one gets
$c = 9072/15692141883605$, whose outside part is $11^{-12}$. Only the
square class matters, because an outside place of even valuation
contributes Hilbert symbol $+1$.
% src: data/l20_admissible.jsonl (normalization block, audited counterexample)

\begin{theorem}[Classical Schinzel H supplies emergent-free members; PROVED implication]
\label{thm:schinzel}
Assume the class certificate (all symbols at $S$, at the moving prime
$Q$, and at infinity equal $+1$ throughout the progression), assume the
square class of $c$ is supported on $S$ and $G(t)$ is a $p$-adic unit for
every $p\in S$, and assume $Q$ and $G$ are irreducible with positive
leading coefficients whose product has no fixed prime divisor.  Then
Schinzel's Hypothesis~H for $\{Q,G\}$ implies that the class contains
infinitely many members with ladder verdict \emph{zero}.
\end{theorem}

\begin{proof}[Proof sketch]
Schinzel supplies infinitely many $t$ with $Q(t)$ and $G(t)$
simultaneously prime; the side conditions (Cauchy bound, $b\ne0,1$, and
finitely many forbidden coincidences) exclude only finitely many $t$.
Put $R=G(t)$.  The part of $c$ outside $S$ need not be trivial, but every
one of its valuations is even.  At the moving prime, $b\equiv0\pmod Q$
and the kernel identity gives
$P(b)\equiv-\delta a^4Z^4A^2\pmod Q$, a unit, so $R\ne Q$.
Consequently $R$ is the unique additional outside place of odd valuation;
every other outside place has even valuation and $v_p(d)=0$, hence symbol
$+1$.  The certificate pins the symbols at $S$, at $Q$, and at infinity
to $+1$.  Hilbert reciprocity forces $(x_0,d_0)_R=+1$.  Thus no place
carries symbol $-1$, and the ladder verdict is \emph{zero}.
% src: data/l18_schinzel_implies_h.jsonl (fixed-class implication)
% src: data/l20_admissible.jsonl (content formula and square-class correction)
% src: THEOREMS.md L9 (reciprocity steering), h10q.py odd-prime symbol formula
\end{proof}

\subsection{Uniform admissibility status}

For the class of Theorem~\ref{thm:classexist}, let
$L(t)=q_1+Nt$ and normalize
$P(fL(t))=cG(t)$ with $c>0$ and $G\in\Z[t]$ primitive.  The exact status
is as follows.

\begin{table}[ht]
\centering
\small
\begin{tabular}{@{}p{.09\textwidth}p{.49\textwidth}p{.30\textwidth}@{}}
\toprule
condition & requirement & uniform status\\
\midrule
(a) & irreducible polynomials with positive leading coefficients &
\textsc{proved} for every fixed $Z\in\Q^\times$; Hilbert selection in the
admissible $a$ progression \textsc{proved}\\
(b) & $\gcd(q_1,N)=1$ & \textsc{proved}\\
(c) & $L(t)G(t)$ has no fixed prime divisor & \textsc{proved}\\
(d) & $G(t)$ is a $p$-adic unit at $S$ and the frozen square classes are
constant & \textsc{proved}\\
\bottomrule
\end{tabular}
\caption{Admissibility of the uniformly constructed square-branch classes.}
\label{tab:admissibility}
\end{table}
% src: data/l20_admissible.jsonl; data/l22_elimination.jsonl (fixed-Z-theorem row)

For (a), the constant and leading coefficients of $P$ are
$P_0=P_8=4a^8AZ^4>0$, and the affine substitution
$b=f(q_1+Nt)$ preserves factorization over $\Q$.  Thus $G$ has positive
lead and
\[
 G\text{ irreducible over }\Q
 \quad\Longleftrightarrow\quad
 P\text{ irreducible over }\Q.
\]
Theorem~\ref{thm:fiber-irred} proves the vertical polynomial irreducible
for every fixed $Z\in\Q^\times$, and
Corollary~\ref{cor:admissible-irred} selects a good $a$ in the required
character progression.  Thus the fixed cell, not merely the generic
two-variable family, is licensed.
% src: data/l20_admissible.jsonl; data/l22_elimination.jsonl

For (b), the prime support of $N$ is exactly $S$, while the CRT base
$q_1$ is a unit modulo every prime in $S$.  For (c), primes in $S$ never
divide either value by the Taylor congruences.  If $p\notin S$, the linear
factor has one root modulo $p$; a fixed divisor would force the degree-$8$
polynomial $G$ to vanish at the other $p-1$ residues, hence $p\le9$.
Every such prime already lies in $S$, a contradiction.  For (d), the
Taylor-exponent construction gives coefficientwise
\[
 G(t)\equiv G(0)\not\equiv0\pmod p
 \quad(p\in S\text{ odd}),\qquad
 G(t)\equiv G(0)\not\equiv0\pmod8
\]
and the analogous congruences for $L(t)$.  The resulting ratios are local
squares, so every frozen Hilbert symbol is constant for all integer $t$.
% src: data/l20_admissible.jsonl (uniform proofs for conditions (a)--(d))

The modular engine remains an exact finite cross-check.  The L20 class
audit certified $353/353$ independently reconstructed grid classes and
$706/706$ legacy cell--branch rows, while the L21 audit certified all
$353/353$ recorded target fibers.  The L22 endpoint--Newton theorem is
strictly stronger in scope: it proves the vertical statement for every
fixed nonzero rational $Z$, without extrapolating from those finite rows.
% src: data/l20_admissible.jsonl; data/l21_irred_generic.jsonl
% src: data/l22_elimination.jsonl (fixed-Z-theorem row)

A separate $24$-row cross-check re-ran the proven-primality kernel end to
end: every sampled member had one proved-prime ladder residual with symbol
$+1$ and verdict \emph{zero}.
% src: data/l18_schinzel_implies_h.jsonl (24 verified_prime_rung rows)

\begin{remark}[Exact scope]
Schinzel supplies members only after a chosen class satisfies every row of
Table~\ref{tab:admissibility}.  Every one of those rows is now PROVED on the
fixed canonical branch: Theorem~\ref{thm:fiber-irred} and
Corollary~\ref{cor:admissible-irred} supply irreducibility, while L20
supplies normalization, positivity, coprimality, absence of a fixed
divisor, and frozen local data.  Therefore the implication has the exact
form
\[
 \text{classical Schinzel H}
 \quad\Longrightarrow\quad
 \text{the conditional six-quantifier conclusion}.
\]
Schinzel~H is the sole conjectural input.  The conclusion remains
conditional: this argument neither proves Schinzel~H nor gives an
unconditional solution of Hilbert's Tenth Problem over $\Q$.
% src: data/l22_elimination.jsonl (fixed-Z-theorem row)
% src: data/l20_admissible.jsonl; data/l18_schinzel_implies_h.jsonl
\end{remark}

\section{The step-(ii) density law and the off-grid frontier}
\label{sec:frontier}

This section reports the two newest components of the program: an
unconditional local bound on the sieve masses together with the exact
decay law it implies, and the current state of the off-grid horizon.

\subsection{An unconditional local bound}

Fix an aligned class with parameters $(\varepsilon,f,q_1,N)$ and write
$F(k) = P(\varepsilon f (q_1 + kN))$ for its degree-$8$ class
polynomial, and $u(k) = 2\alpha\varepsilon f(q_1+kN)$ for the unit whose
Legendre symbol decides bad signature. For an odd prime $p$ let $m_p$
denote the Haar mass in $\Zp$ of the set of $k$ with $v_p(F(k))$ odd and
$\leg{u(k)}{p} = -1$: the local mass of the emergent obstruction.

\begin{theorem}[Local sieve bound; PROVED]
\label{thm:localbound}
For every aligned class of the grid,
\[
   m_3 = m_5 = m_7 = 0,
   \qquad\text{and}\qquad
   m_p \;\le\; \frac{8}{p+1} \;<\; 1
   \quad\text{for all } p \ge 11 ,
\]
the bound holding for singular roots as well, counted with
multiplicity.
% src: data/l17_stepii.jsonl (meta.labels.theorem_replays = PROVED)
\end{theorem}

The vanishing at $p \in \{3,5,7\}$ is exactly the frozen-alignment
condition of the class certificate: those places are controlled, and
their symbols are pinned to $+1$ by construction. For $p \ge 11$ the
bound is the root count of a degree-$8$ polynomial: at most eight
residues can be roots, each simple root contributes Haar mass
$1/(p+1)$ for odd valuation, and a singular root contributes at most
its multiplicity times that mass, which is again bounded by the total
degree. The bound was verified independently on every stored local
pair: $102{,}439$ pairs with $p \le 10^4$ across the $293$ classes,
zero violations.
% src: data/l17_stepii.jsonl (verification.positive_local_pairs_checked_p_le_10000 = 102439, bound_violations = 0)

Theorem~\ref{thm:localbound} settles a natural hope in the negative,
and this is the important structural point.

\begin{corollary}[No uniform positive mass; PROVED for the current classes]
\label{cor:nouniform}
For each of the $293$ classes the splitting condition governing the
bad-signature event has Chebotarev exponent $c = 1/2$: the class
polynomial is irreducible with a replayed degree-$8$ certificate, and a
replayed witness shows the bad unit is a nonsquare at a simple bad
root. Consequently the truncated clean product obeys
\[
   \prod_{p \le X} (1 - m_p) \;\asymp\; (\log X)^{-1/2},
\]
so no cutoff-uniform positive lower bound on the clean mass exists:
the small-layer clean probability decays, slowly, to zero.
% src: data/l17_stepii.jsonl (verdict; serialized_witnesses_replayed = 293)
\end{corollary}

The measured exponent agrees: the root-harmonic estimate gives
$c \approx 0.4778$, the product estimate $c \approx 0.4801$, and the
same-member sieve between $10^4$ and $10^5$ gives $c \approx 0.4590$
(EVIDENCE).
% src: data/l17_stepii.jsonl (evidence.root_harmonic_c_mean, product_c_mean, same_member_sieve_c_10000_100000)

The consequence for hypothesis $H$ is a sharpening, not a defeat. The
number of admissible prime members up to height $X$ in a class grows
like $X/(\log X)^{2}$ by the Dirichlet--Bateman--Horn heuristic
\cite{batemanhorn}, while
each carries clean probability $\asymp (\log X)^{-1/2}$; the expected
count of emergent-free members therefore still diverges, at rate
$X/(\log X)^{5/2}$. What Corollary~\ref{cor:nouniform} rules out is the
clean shortcut --- a uniform positive constant --- and what it leaves is
a divergence statement of exactly the shape the measured counting law
already exhibits (first emergent-free member at $k=0$ in $207$ of $293$
classes).

\subsection{The off-grid horizon}

The grid is $w$ prime, $w<100$.  Off-grid closures now exist for
\[
 w\in\{101,103,107,109,113,127,131,137,139,149,151,157,163,167,173,179\},
 \qquad u=-1,
\]
together with $[131,[5,1]]$, each certified by an aligned class and an exact
ladder verdict \emph{zero}.  The routes are the L11 family, the canonical
fixed-factor escape, and the non-canonical $a=7$, $\tau=0$ family.  The
last searched open cell, $[179,[-1,1]]$, is now closed, so every off-grid
cell tried to date has a replayed member.  This is a finite-horizon
statement only.
% src: data/l17_horizon*.jsonl; data/l18_horizon_sweep_a.jsonl
% src: data/l18_horizon_sweep_b.jsonl; data/l18_route131.jsonl
% src: data/l19_cell179.jsonl (closure and summary rows)

Two negative or delimiting findings carry structural content.

\begin{theorem}[The canonical $5$-wall; PROVED]
\label{thm:fivewall}
On the canonical escape route ($z=-w$, $a=1$, $\tau=3/5$, $f=w$), write
$e = v_5(D)$. Then $v_5(x_0) = -3-2e$ is odd and $v_5(d_0)=0$, so
\[
   \mathrm{Hilb}_5 \;=\; -\leg{w}{5}\leg{q_1}{5},
   \qquad
   \mathrm{Hilb}_{q_1} \;=\; \leg{5}{q_1} ,
\]
and both equal $+1$ for some admissible $q_1$ if and only if
$\leg{w}{5} = -1$. Since the L11 route requires a root of
$1+4r^2 \equiv 0 \pmod w$, i.e.\ $w \equiv 1 \pmod 4$, the canonical
protocol is obstructed exactly for
\[
   w \equiv 3 \pmod 4 \ \text{ and } \ \leg{w}{5} = +1
   \qquad\Longleftrightarrow\qquad
   w \equiv 11, 19 \pmod{20}.
\]
% src: data/l18_fivewall.jsonl (proved identity + 53-prime audit, 0 mismatches)
\end{theorem}

The criterion was audited on all $53$ primes $w \in [101,400]$ with both
character branches of $q_1$ covered: $26$ cells have the L11 route, $13$
the canonical escape route, $14$ are obstructed, and there were zero
identity or classification mismatches.
% src: data/l18_fivewall.jsonl (classification_counts, identity_mismatches = 0)

The wall is a statement about the canonical parameter, not about the cells.
In fact it is always breakable on the class side.

\begin{theorem}[Every canonical $5$-wall is breakable; PROVED]
\label{prop:break}
Let $w\equiv11,19\pmod{20}$ be prime and set $z=-w$, $f=w$.  There is an
odd integer $a$ such that the $\tau=0$ fixed-factor family has an aligned
class for infinitely many positive primes $q_1$.  This theorem concerns
class existence; it does not assert an emergent-free member.  If $a$ is
fixed to $7$, the class-side break exists exactly when
$\leg{w}{197}=-1$.
\end{theorem}
% src: data/l19_tauzero.jsonl (criterion and always-breakable theorem rows)

\begin{proof}
Because $w\equiv3\pmod4$, the character sum
$\sum_{r\bmod w}\leg{1+4r^2}{w}=-1$ has no zero terms.  Exactly
$(w+1)/2$ residues therefore give character $-1$; choose one for $a\bmod
w$.  By CRT also impose
$a\equiv0\pmod{\operatorname{rad}(w^3+2)}$ and choose an odd lift.  For
$A=1+4a^2$, $Z=-w^3$, and $D=1-Z-a^2Z^2$, the identity
\[
 4D\equiv(Z-2)^2\pmod p\qquad(p\mid A)
\]
then gives $\gcd(A,D)=1$ while preserving $\leg{A}{w}=-1$.

At every prime occurring to odd order in $A$, prescribe
$\leg{-2\varepsilon wq_1}{p}=+1$; impose the exact $2$-adic class
$\varepsilon wq_1\equiv a\pmod4$; and make
$-2A\varepsilon wq_1$ a square at each of $3,5,7$ not dividing $A$.
These independent local conditions define a reduced CRT class, so
Dirichlet gives infinitely many primes $q_1$ in it.  After finitely many
exclusions and a real-sign threshold, every frozen, moving, fixed-factor,
and infinite symbol is $+1$, and the exponent lemma freezes the resulting
progression.  The factor-by-factor identity
\[
 \prod_{p\mid A}(x,d)_p\,\leg{A}{q_1}
   =\leg{-2\varepsilon}{A}\leg{w}{A}
\]
shows that for odd $a$ the break criterion is $\leg{w}{A}=-1$.
Putting $a=7$ gives $A=197$ and the stated fixed-parameter criterion.
\end{proof}
% src: data/l19_tauzero.jsonl (algebraic-criterion and uniform-proof rows)

\begin{proposition}[Finite member closures; PROVED per row]
The non-canonical $\tau=0$ family closes the obstructed cells
$w=131,139,151$ at $a=7$, $q_1=Q=41$, $k=0$.  It also closes
$[179,[-1,1]]$ at $a=7$, $\varepsilon=-1$, $f=179$,
$q_1=Q=251$, $k=0$, with prime rung, tied status exact, and empty
ramified set.
\end{proposition}
% src: data/l18_route131.jsonl; data/l18_horizon_sweep_a.jsonl
% src: data/l19_cell179.jsonl (verified closure row)

These member rows are stronger than the class theorem only at their four
individual cells.  Their fixed factor $f=w$ divides $\num(z)$, so they pass
through the controlled non-coprimality door excluded by the generalized
coprime wall.  The kernel verdict uses the same tied-status, ramification,
alignment, and complete cofactor ladder as the canonical rows; refusals in
the surrounding scans are not evidence.
% src: data/l18_route131.jsonl; data/l18_horizon_sweep_a.jsonl
% src: data/l19_cell179.jsonl

\begin{remark}[Composite $w$]
Composite $w$ is not admissible in the architecture: $w$ indexes a
place, and the residue-field construction has no meaning for composite
modulus --- the break occurs already at the field constructor. The
downstream arithmetic does generalize once $w$ is decomposed into
genuine prime places; as an illustration the pseudo-cell
$[21,[-1,1]]$ closes with the escape route at the honest prime $f = 3$
(member $k=0$, $Q=11$, factorint rung, tied and ramified checks exact).
This is a clarification of scope, not an extension of the horizon.
% src: data/l17_composite_w.jsonl
\end{remark}


\section{Conclusion and outlook}

The program's current position is sharp. The algebraic architecture is
complete and machine-checked; the search space is delimited by two
theorems (the generalized wall, valid for composite squarefree $A$, and
the class characterization); the finite verification is exhaustive on the
grid and extends off-grid to $w \in \{101,103,107,109,113,127,137\}$; and the
residual analytic content has been isolated to a single classical
statement about prime values of a degree-$8$ polynomial in an arithmetic
progression, with all of its local conditions verified.

Three directions carry the most weight.

\begin{enumerate}
\item \emph{The step-(ii) divergence.} Theorem~\ref{thm:localbound} and
Corollary~\ref{cor:nouniform} settle the local question: the masses are
bounded by $8/(p+1)$, the frozen primes contribute nothing, and the
clean product decays like $(\log X)^{-1/2}$, so no uniform positive
mass exists. What remains is to convert the resulting divergence of the
expected count of emergent-free members, at rate $X/(\log X)^{5/2}$,
into an existence statement; that is precisely where the degree-$8$
prime-value input enters.
\item \emph{The protocol wall, answered.} Composite $w$ is out of scope
by construction ($w$ indexes a place). The question raised by the first
obstructed cell --- for which primes $w$ does the canonical escape route
collide with the L11 route at the frozen place $5$? --- is now answered
exactly by Theorem~\ref{thm:fivewall}: $w \equiv 11, 19 \pmod{20}$. It is
also no longer an obstruction, since Proposition~\ref{prop:break} closes
every such cell met so far off the canonical protocol. No off-grid cell
is currently known to be intrinsically obstructed, and the natural next
question is whether the $a = 7$, $\tau = 0$ family closes the obstructed
congruence class in general.
\item \emph{Degree-$8$ prime values.} Unconditional prime-value theorems
stop at degree three. Any unconditional progress for a single grid class
would convert the conditional statement into an unconditional one for
that cell, and the audit shows there is no local obstruction anywhere on
the grid.
\end{enumerate}

\section*{Reproducibility}

All certificates in this paper are produced by a stdlib-only kernel with
a proven-primality discipline: deterministic Miller--Rabin below
$3.3\times 10^{24}$, Pocklington with the Brillhart--Lehmer--Selfridge
relaxation above, and explicit refusal labels that never count as
evidence. The companion paper documents the engine, the artifact
schemas, the exact reproduction commands, and the precise scope of what
is regenerable from checked-in scripts.

\begin{thebibliography}{99}

\bibitem{batemanhorn}
P.~T. Bateman and R.~A. Horn,
\emph{A heuristic asymptotic formula concerning the distribution of prime
numbers},
Mathematics of Computation \textbf{16} (1962), 363--367.

\bibitem{ctsansuc}
J.-L. Colliot-Th\'el\`ene and J.-J. Sansuc,
\emph{Sur le principe de Hasse et l'approximation faible, et sur une
hypoth\`ese de Schinzel},
Acta Arithmetica \textbf{41} (1982), no.~1, 33--53.

\bibitem{ctssd}
J.-L. Colliot-Th\'el\`ene, A.~N. Skorobogatov and Sir P. Swinnerton-Dyer,
\emph{Rational points and zero-cycles on fibred varieties: Schinzel's
hypothesis and Salberger's device},
Journal f\"ur die reine und angewandte Mathematik \textbf{495} (1998),
1--28.

\bibitem{cz}
G.~Cornelissen and K.~Zahidi,
\emph{Topology of Diophantine sets: remarks on Mazur's conjectures},
in: Hilbert's Tenth Problem: Relations with Arithmetic and Algebraic
Geometry (Ghent, 1999), Contemporary Mathematics \textbf{270},
American Mathematical Society, Providence RI, 2000, pp.~253--260;
arXiv:math/0006140.

\bibitem{daans}
N.~Daans,
\emph{Universally defining $\Z$ in $\Q$ with $10$ quantifiers},
Journal of the London Mathematical Society (2) \textbf{109} (2024),
no.~2, e12864, 20 pp.;
\texttt{doi:10.1112/jlms.12864}; arXiv:2301.02107.

\bibitem{ddf}
N.~Daans, P.~Dittmann and A.~Fehm,
\emph{Existential rank and essential dimension of diophantine sets},
preprint, arXiv:2102.06941v5 (2021--2024).

\bibitem{hsw}
Y.~Harpaz, A.~N. Skorobogatov and O.~Wittenberg,
\emph{The Hardy--Littlewood conjecture and rational points},
Compositio Mathematica \textbf{150} (2014), no.~12, 2095--2111;
\texttt{doi:10.1112/S0010437X14007568}; arXiv:1304.3333.

\bibitem{helfgott}
H.~A. Helfgott,
\emph{Square-free values of $f(p)$, $f$ cubic},
Acta Mathematica \textbf{213} (2014), no.~1, 107--135;
\texttt{doi:10.1007/s11511-014-0117-2}; arXiv:1112.3820.

\bibitem{koenigsmann}
J.~Koenigsmann,
\emph{Defining $\Z$ in $\Q$},
Annals of Mathematics (2) \textbf{183} (2016), no.~1, 73--93;
\texttt{doi:10.4007/annals.2016.183.1.2}; arXiv:1011.3424.

\bibitem{krumm}
D.~Krumm,
\emph{Squarefree parts of polynomial values},
Journal de Th\'eorie des Nombres de Bordeaux \textbf{28} (2016), no.~3,
699--724;
\texttt{doi:10.5802/jtnb.959}; arXiv:1407.4890.

\bibitem{reuss}
T.~Reuss,
\emph{Power-free values of polynomials},
Bulletin of the London Mathematical Society \textbf{47} (2015), no.~2,
270--284;
\texttt{doi:10.1112/blms/bdv007}; arXiv:1307.2802.

\bibitem{schinzel}
A.~Schinzel and W.~Sierpi\'nski,
\emph{Sur certaines hypoth\`eses concernant les nombres premiers},
Acta Arithmetica \textbf{4} (1958), 185--208;
erratum, ibid.\ \textbf{5} (1958), 259.

\bibitem{sun}
Z.-W. Sun,
\emph{$\Q\setminus\Z$ is diophantine over $\Q$ with $7$ unknowns},
preprint, arXiv:2607.28606 (unrefereed).

\end{thebibliography}

\end{document}
